% Chapter 1 -- Introduction (\input by main.tex)
% Frames the problem, the modular approach, the two research questions, the
% contribution, the scope, and the thesis outline.
% Published literature uses BibTeX \cite{}.
%
% CONSISTENCY CONSTRAINTS (checked at draft time):
%   - Every claim here is a subset of the claims map (C1--C11).
%     Nothing is stated beyond the PARTIALLY-HOLDS scope of the held-out completion
%     result (C9). The single scientific contribution (the donor metric) matches
%     Chapter 2's Sec 2.5; the two diagnostic findings (recall root-cause /
%     synthetic-pretraining redundancy) are reported as findings, not contributions;
%     the pipeline, runtime characterisation, and amodal-GT builder are named as
%     ENGINEERING contributions, and tracking is explicitly not a contribution.
%
% Honesty-phrasing rules applied (personal_agenda P1):
%   - recall is stated against its car-only, >=10-surviving-points denominator
%   - the held-out completion result is PARTIALLY HOLDS, never unqualified
%   - negative results are named as findings, not buried
%   - the detection headline is a single-sequence (seq 08) result, stated as such
%
\chapter{Introduction}
\label{ch:introduction}

\section{Background and Rationale}
\label{sec:intro-background}

Autonomous driving depends on a perception system that can recover, several times a
second, the three-dimensional layout of the scene around the vehicle: where the other
road users are, how large they are, and which way they face. LiDAR is central to this
task. It measures range directly, returning an unordered cloud of over a hundred
thousand points per scan, largely invariant to lighting and carrying the metric depth
that motion planning needs. The benchmarks that opened the field were built around it:
the KITTI suite~\cite{geiger2012we} paired a scanning LiDAR with calibrated cameras,
and its point-wise successor SemanticKITTI~\cite{Behley_2019_ICCV} turned the same
sequences into the dense segmentation benchmark much of the field is now measured
against.

The dominant paradigm is to train a deep network end to end on annotated scans:
point-based segmentation networks~\cite{Qi_2017_CVPR,NIPS2017_d8bf84be,Hu_2020_CVPR},
projection- and voxel-based models~\cite{10.1109/ICRA.2018.8462926,milioto2019rangenet++,zhu2021cylindrical},
single-pass object detectors~\cite{lang2019pointpillars,yin2021centerpoint}, and
panoptic models~\cite{zhou2021panopticpolarnet,hong2021dsnet}. These define the state
of the art in accuracy, which is paid for in two costs. The first is supervision: the
results rest on dense, point-wise labels, and the annotation behind a benchmark like
SemanticKITTI --- a class on every one of billions of points --- is out of reach for
most groups and most new datasets. The second is opacity: an end-to-end network spreads
each decision across millions of weights, so when it misses a car there is no single
stage to inspect or repair.

This thesis takes the opposite starting point, asking how far a \emph{modular},
largely annotation-free pipeline can go: classical geometry for the parts geometry
solves well, and small learned models only where discrimination needs them --- an
approach to LiDAR vehicle detection that has drawn renewed interest~\cite{lim2025longrange}.
Clustering proposes candidate objects, a small binary \emph{car}/\emph{not-car}
classifier keeps the vehicles, and each detected car is completed by a PCN network;
every stage is separately inspectable, and the learned components are trained without
point-wise labels. Chapter~\ref{ch:pipeline} gives the full pipeline
(Figure~\ref{fig:pipeline}).

\section{Challenges in Real-Data Completion}
\label{sec:intro-challenges}

It is the second half of this pipeline that raises the problem the thesis is built
around. Completing a partially observed car --- recovering the far side the sensor
cannot reach --- is attractive for downstream box estimation, and networks such as
PCN~\cite{yuan2019pcnpointcompletionnetwork} are well studied on synthetic benchmarks.
On real LiDAR they cannot be evaluated: there is no clean ground-truth complete shape,
and the obvious substitute --- accumulating the same static car across many frames ---
is itself one-sided and, as this thesis shows, invalid as a reference, because the raw
partial input scores best against it. Before completion can be claimed to help,
it must be possible to measure at all. This is where the topic narrowed: an initial aim
of applying learned completion to real LiDAR ran into the fact that its quality could
not be measured validly at all, and that gap --- not the completion network itself ---
became the question the thesis is built around.

\section{Research Objectives}
\label{sec:intro-objectives}

The work is organised around two research questions about measuring and validating
completion on real LiDAR; the central contribution answers them, and two further
findings --- reported below in their own right --- record what the investigation
revealed about the detection front-end along the way.

\begin{enumerate}
  \item \textbf{A Valid Evaluation Method for Real-LiDAR Completion.}
    Establish how the quality of point-cloud completion can be evaluated on
    \emph{real} LiDAR, where no clean reference shape exists and an accumulated
    pseudo-ground-truth is itself one-sided --- a measurement question prior to any
    claim that completion helps (RQ1; answered in
    Section~\ref{sec:donor-validation}).
  \item \textbf{Evidence That Completion Recovers Unseen Surface.}
    Determine whether learned completion recovers genuinely \emph{unseen} surface,
    as opposed to redistributing the points it was already given; whether it
    improves downstream amodal bounding boxes on real LiDAR cars; and whether any
    such benefit survives on a held-out sequence under a pre-registered test (RQ2;
    answered in Section~\ref{sec:donor-validation}, with the held-out outcome
    reported exactly as it fell).
\end{enumerate}

\section{Research Contributions}
\label{sec:intro-contributions}

The thesis makes a single scientific contribution, methodological rather than
architectural, together with the engineering contributions that make it possible.

\begin{itemize}
  \item \textbf{A Validated Metric for Occluded-Vehicle Completion on Real LiDAR.}
    Chamfer distance against an accumulated pseudo-ground-truth is invalid on real
    cars, so the thesis defines a donor-frame occluded-side coverage metric: it
    scores how much of the surface seen only in other frames of the same
    static car a completion recovers, with a hallucination guard against surface
    invented outside the car's box, and validates it against a four-item gate.
    Applied to real sequence-08 cars, it shows that PCN recovers occluded surface
    several times above the symmetry-mirror baseline and improves amodal boxes on
    static, gate-passed cars, and that the benefit partially holds on a
    pre-registered held-out sequence (Section~\ref{sec:donor-validation}). The
    detection and completion stages are established building blocks, reimplemented
    and tuned here; the contribution is the measurement built on top of them.

  \item \textbf{A Modular, Near-Annotation-Free Detection-and-Completion Pipeline.}
    A pipeline in which classical geometry does the work geometry solves well and
    small learned models are used only where discrimination needs them, with every
    stage separately inspectable and its per-stage runtime cost measured rather
    than estimated.

  \item \textbf{An Amodal Ground-Truth Box Builder.}
    A procedure that accumulates a static car's frames into an amodal
    bounding-box reference, used to measure completion's effect on downstream box
    estimation where no external amodal reference exists.

  \item \textbf{Reproducibility Infrastructure.}
    The frozen configuration, deterministic seeds, and evaluation scripts on which
    every reported number rests.
\end{itemize}

Along the way the investigation produced two diagnostic findings: that recall is
limited by density clustering splitting single cars into fragments rather than by the
classifier, with a chain of repair strategies that fail to generalise
(Section~\ref{sec:recall-analysis}); and that synthetic pretraining is redundant for
the classifier once real training clusters are plentiful (Chapter~\ref{ch:evaluation}).

The directions that did not work (architecture swaps, real-data completion
fine-tuning, a sparse-fallback length model) are reported as negative findings
rather than omitted.

\textbf{The detector is car-only:} the classifier is binary (car / not-car). \textbf{The tracker is a filter only:} it removes flicker
detections and is not counted among the contributions. \textbf{Completion is trained
without real labels:} the network sees only synthetic partial renders, never a labelled
real car --- what makes the pipeline near-annotation-free, and the source of the
sim-to-real questions Chapter~\ref{ch:evaluation} takes up. \textbf{The system is
offline:} it runs at about \SI{623.5}{\milli\second} per frame (\num{1.6}~fps), and
real-time is out of scope (Section~\ref{sec:results-runtime}). \textbf{Detection is
reported at one operating point but cross-validated:} sequence~08 (\num{4071} scans) is
where the ablations, recall root-cause, and completion analysis are decomposed, and a
leave-one-sequence-out cross-validation (Section~\ref{sec:eval-loso}) bounds the
model-selection exposure to about two precision points. Throughout, ``recall'' is
reported against its car-only, at-least-ten-surviving-points denominator
(Section~\ref{sec:eval-protocol}).

\textbf{Approach and method.} The investigation is empirical and proceeds in four
steps, each held to a single frozen configuration for reproducibility. First, we
assemble the modular pipeline and train its learned components without point-wise
labels: ground removal and density clustering propose candidate objects, a binary
classifier keeps the cars, and PCN completes each one. Second, we fix a reproducible
detection protocol --- point-level IoU matching against SemanticKITTI-derived
ground truth, scored by precision, recall, F\textsubscript{1}, and mean IoU --- and
characterise the detector on real KITTI-style sequences, with leave-one-sequence-out
cross-validation guarding against model-selection bias. Third, because Chamfer distance
against an accumulated pseudo-ground-truth is invalid on real cars, we define the
donor-frame coverage metric and put it through a four-item validation gate. Fourth, we
use the validated metric to measure whether completion recovers genuinely unseen surface
and improves amodal boxes, with the outcome reported exactly as it fell on a
pre-registered held-out sequence. The perspective throughout is that geometry should be left to classical
methods and learning reserved for where discrimination needs it, and that a completion
claim on real data is only as trustworthy as the measurement behind it.

\section{Thesis Organization}
\label{sec:intro-outline}

The remainder of the thesis is organised into five chapters.

\textbf{Chapter~\ref{ch:introduction}: Introduction} presents the background and
rationale, the challenges of evaluating completion on real LiDAR, the research
objectives, the contributions, and the organisation of the thesis.

\textbf{Chapter~\ref{ch:background}: Literature Review \& Theoretical Background}
places the pipeline against the deep-learning segmentation and point-completion
literature, carries the positioning table that defends the comparison, and fixes
the notation and standard definitions for the operators the pipeline is built from.

\textbf{Chapter~\ref{ch:pipeline}: Proposed Methodology} describes every stage of
the detection pipeline and the classifier's training, the point-cloud completion
method, and --- as the chapter's destination --- the definition of the donor-frame
coverage metric that is this thesis's contribution.

\textbf{Chapter~\ref{ch:evaluation}: Experimental Evaluation} evaluates the pipeline
end to end: it fixes the detection protocol and reports the seq-08 operating point,
the ablations, the distance breakdown, and the runtime; root-causes the recall
shortfall to clustering fragmentation and reports the chain of negative repair
strategies, which is also the bridge to completion; and turns to its centre, putting
the donor metric through its validation battery and using it to report the real-data
coverage results, the amodal-box utility, the moving-car case, and the pre-registered
held-out replication.

\textbf{Chapter~\ref{ch:conclusion}: Conclusion and Future Work} answers the two
research questions, names every weakness the work is subject to, and lays out what
remains.
